\section{Conclusion}\label{sec:conclusion}

\subsection{Summary and Discussion}
\label{subsec:conclusion-summary}

This project worked through a sequence of increasingly complex quantum systems: from single-qubit operations and Bell state entanglement, through two-level and two-qubit Hamiltonians, to the Lipkin--Meshkov--Glick model for $N=2$ and $N=4$ particles. At each stage, classical diagonalization provided exact reference solutions against which the VQE results were compared.

For the two-level Hamiltonian the variational approach performs essentially perfectly: the single-parameter $R_y$ ansatz is exactly expressive, and errors below $10^{-10}$ across the full $\lambda$ range confirm the ansatz is not the limiting factor. The two-qubit case is more revealing. The $R_y$-CNOT-$R_y$ circuit is also exactly expressive, and for $\lambda\lesssim0.4$ it matches the exact ground state to machine precision. At the avoided crossing near $\lambda\approx0.4$, however, both VQE implementations diverge from the exact solution: the optimizer fails to follow the ground state to the lower branch and instead tracks the first excited state, with errors jumping from $\sim10^{-11}$ to $\sim10^{-1}$ at the crossing. This is a branch-tracking failure rather than an ansatz limitation — the circuit is capable of representing the correct state, but cold-start gradient descent has no mechanism to switch branches. The entanglement analysis was clean up to the crossing: the von Neumann entropy correctly distinguishes product states from maximally entangled Bell states, and in the two-qubit Hamiltonian it rises sharply near the level repulsion, where the ground state becomes a strongly correlated superposition. This connection between level repulsion and entanglement growth is one of the cleaner results in the project.

The Lipkin model is where things get harder. Classical diagonalization of both the $J=1$ and $J=2$ Hamiltonians works without difficulty, and the Hartree-Fock comparison confirms the expected behavior: exact below the critical coupling, systematically biased above it. The VQE also handles $J=1$ well. For $J=2$ the residual error of $\approx10^{-2}$ reflects a genuine structural limitation: the four-qubit Hilbert space contains three symmetry sectors, and the hardware-efficient ansatz has no mechanism to stay within the physical one. The problem is worst near the phase transition, where the optimization landscape is locally flat and gradient-based methods stall. Four random restarts help but do not fix it. These results are also obtained under ideal conditions — exact statevector simulation, no gate noise. On real hardware the precision would be worse, likely by a significant margin at circuit depth 4.

\subsection{Reflections}
\label{subsec:conclusion-reflections}

Working through this progression made the formalism more concrete in ways that are hard to get from reading alone. The link between entanglement and interaction strength — visible in the entropy curves as $\lambda$ passes through the crossing region — and the way the Pauli decomposition translates a physical Hamiltonian into a form that can be measured circuit-by-circuit both became much clearer through the implementation. The VQE errors were instructive in a way I did not fully anticipate. For the two-level system, cold-start BFGS converges immediately and perfectly. For the two-qubit system, the optimizer succeeds up to the avoided crossing, then fails in a precisely located and physically interpretable way — the crossing point where the ground state switches character. For the Lipkin $J=2$ case the failure is again tied to the physics, this time the near-flat landscape near the phase transition rather than a branch discontinuity. In all three cases the optimization difficulty maps directly onto the structure of the problem, which I found more interesting than a generic convergence failure would have been.

Having two independent VQE implementations — manual and Qiskit-based — agree throughout was also reassuring: to machine precision for $\lambda\lesssim0.4$ in the two-qubit case, and identically in failure for $\lambda>0.4$. The latter is particularly useful — both implementations making the same mistake at the same crossing point confirms the behavior is physical, not a software artifact, and ruled out implementation errors before moving to the more demanding Lipkin system.

\subsection{Future Work}
\label{subsec:conclusion-future-work}

For the two-qubit system, the branch-tracking failure at $\lambda\approx0.4$ could be addressed by warm-starting each $\lambda$ step from the previous solution and supplementing with random restarts near the crossing, or by using an excited-state VQE variant that explicitly targets the lower branch. For the $J=2$ Lipkin system, the most direct improvement would be a symmetry-adapted VQE using a circuit constrained to the $J=2$ sector from the start. This would likely close the $\approx10^{-2}$ precision gap and scale more favorably to larger $N$. A natural extension from there is a sweep over both $V$ and $W$, mapping the full two-parameter phase diagram of the LMG model — the current implementation supports arbitrary $W$, so this is mostly a matter of running the experiments. Finally, testing these circuits on actual quantum hardware via the IBM Quantum network would expose the gap between statevector simulation and real device performance, and would make the limitations identified here more concrete.
